\subsection{Mathematical Research and Integrity}

Mathematical research cycles through discovery, attempted refutation, proof, verification, and communication rather than proceeding linearly from question to theorem.

\begin{process} \label{proc:mathematical-research-cycle} \hlt{Mathematical Research Cycle}\\
For a proposed mathematical problem, maintain the following revisable cycle.
\begin{enumerate}[label=(\roman*)]
\setlength{\itemsep}{0pt}
\item \hlt{Question}: state the objects, scope, current assumptions, and a precise target that can be proved or refuted.
\item \hlt{Literature}: record definitions, earlier results, counterexamples, nearby formulations, and current open status.
\item \hlt{Exploration}: compute examples, boundary cases, and nonexamples while recording the tested domain.
\item \hlt{Conjecture}: state the strongest precise claim supported by exploration and include every suspected hypothesis.
\item \hlt{Attack}: search independently for a proof and a counterexample; record failed routes and new obstructions.
\item \hlt{Revision}: weaken, strengthen, split, or reframe the conjecture when an example or proof obligation demands it.
\item \hlt{Verification}: check every proof dependency, compare with prior results, and seek independent criticism.
\item \hlt{Communication}: state novelty accurately, write a complete proof, cite provenance, and record limitations.
\end{enumerate}
MIT's project laboratory similarly presents mathematical discovery through data generation, conjecture, experiment, proof or counterexample, and writing \citep{mit2013projectlaboratory}.
\end{process}

\begin{process} \label{proc:research-claim-ledger} \hlt{Research Claim Ledger}\\
Maintain one auditable row per substantive claim.
\begin{enumerate}[label=(\roman*)]
\setlength{\itemsep}{0pt}
\item \hlt{Claim}: record the exact wording and scope.
\item \hlt{Status}: record every applicable status: proved, conjectured, computed, observed, reproduced, replicated, or interpreted.
\item \hlt{Support}: record its evidence, dependencies, source locator, and unresolved checks.
\item \hlt{Revision}: update the row whenever a hypothesis, dataset, code revision, or conclusion changes.
\end{enumerate}
The ledger prevents an exploratory observation from silently becoming a theorem or empirical conclusion.
\end{process}

\begin{convention} \label{conv:reproducibility-replicability-generalizability} \hlt{Reproducibility, Replicability, and Generalizability}\\
Adopt the National Academies' 2019 terminology \citep{nasem2019reproducibility}.
\begin{enumerate}[label=(\roman*)]
\setlength{\itemsep}{0pt}
\item \hlt{Reproducibility}: consistent computational results from the same data, code, methods, and analysis conditions.
\item \hlt{Replicability}: consistent results from a new study addressing the same question with newly obtained data.
\item \hlt{Generalizability}: applicability of a result to contexts or populations beyond those directly studied.
\end{enumerate}
These are distinct evidential targets: a reproducible computation need not replicate on new data or generalize to a new setting.
\end{convention}

\begin{principle} \label{princ:research-integrity-record} \hlt{Research Integrity and Stewardship}\\
Maintain research records that permit scrutiny, correction, and appropriate reuse. Report methods and findings honestly; preserve accurate records; credit contributions and prior work; disclose conflicts; respect confidentiality, consent, licenses, and governing rules; and correct the record when an error is found. These duties reflect the honesty, accountability, fairness, and stewardship principles of the Singapore Statement \citep{singapore2010integrity}.
\end{principle}

\begin{warning} \label{warn:novelty-and-negative-results} \hlt{Novelty and Negative Results}\\
Novelty does not imply truth, and publication does not certify correctness. Search for prior and concurrent work before claiming priority, and report a failed proof, null result, or nonreplication with its actual protocol and uncertainty. A negative outcome may diagnose a false conjecture, an inadequate method, or a boundary of generality; the evidence must distinguish these possibilities.
\end{warning}
